\section{Arquitectura técnica}
La solución se define sobre una arquitectura modular que separa interfaz, orquestación de servidor y servicios especializados por responsabilidad.
La descomposición de esta arquitectura puede leerse junto con la Figura~\ref{fig:paquetes-sistema}, que organiza las dependencias entre interfaz, núcleo, servicios, módulos y datos.

\section{Patrones de diseño aplicados}
\begin{itemize}
  \item Separación de responsabilidades por capas y servicios.
  \item Modularización de funcionalidades de alto acoplamiento.
  \item Construcción reactiva para actualización incremental de vistas.
  \item Normalización de utilidades transversales (filtros, paginación, exportación).
\end{itemize}

\section{Diseño de componentes}
\subsection{Interfaz}
Componentes visuales para autenticación, navegación, filtros, resultados y control administrativo.

\subsection{Servidor}
Coordinación de estado de sesión, renderizado dinámico y llamadas a servicios.

\subsection{Servicios}
Funciones de negocio para autenticación, filtrado documental, taxonomías, personas, auditoría y exportación.

\section{Escalabilidad técnica}
\begin{itemize}
  \item Evolucion a persistencia relacional.
  \item Desacople de almacenamiento de activos digitales.
  \item Integración progresiva de pipelines de analítica y búsqueda semántica.
\end{itemize}
La evolución técnica descrita aquí se alinea con la Figura~\ref{fig:gobernanza-rbac} para control de acceso y con la Figura~\ref{fig:despliegue-continuidad} para el ciclo operativo.

\section{Observabilidad técnica}
Definir en versión final un esquema de monitoreo de eventos, métricas de uso, errores de aplicación y tiempos de respuesta por flujo crítico.
